% =====================================================================
\section{Background and related work}
\label{sec:related}
% Facts about cited papers come from the reading-list extractions in
% ../Reading List/Tier_A/Tier_A_extraction.md and Tier_B/*.md, which were
% read against the PDFs. Every number quoted here carries the denominator or
% caveat those extractions flag. The Tier-C papers (JOURNAL_PLAN section 2)
% were read in full on 2026-10-20 before they were cited: AgentSpec and
% Reddy et al. (Section 2.1), AgenTracer and Who&When (Section 2.3), and
% Defects4J (Section 2.4).
% =====================================================================

\subsection{Verifying LLM pipeline outputs at run time}
\label{sec:related-rv}

Runtime verification checks behaviour against a specification during
execution~\citep{leucker2009rv}. Multi-agent frameworks divide tasks among
specialised agents \citep{metagpt,camel,chatdev,autogen}; coordination remains
an open problem~\citep{guo2024survey}. Output schemas do not establish that a
stage's claims are grounded in its inputs. LLM judges \citep{geval,mtbench}
can assess outputs but exhibit position and self-preference biases
\citep{positionbias,selfpreference}. External-tool feedback can improve
correction but needs a trusted resource~\citep{critic}. Our deterministic
checks use pipeline state and detect missing provenance, not falsity
\citep{hoang2026eas}. They also identify the check class and location by
construction, a difficult post-hoc diagnosis problem
\citep{tan2026agentchaos}.

AgentSpec enforces event-triggered policies~\citep{wang2026agentspec};
\citet{reddy2026gates} use deterministic gates before an airline agent's tool
calls, raising task success from 29.6\% to 42.0\% with gpt-4o-mini. Those
systems check actions against policies. Our supervisor checks stage outputs
against upstream state, and compares informative and generic rejections.

\subsection{Self-correction and the role of feedback}
\label{sec:related-selfcorrection}

A large literature studies whether a model can improve an output when told,
or when it tells itself, that the output is wrong. In Self-Refine the model
critiques and revises its own output~\citep{selfrefine}, and Reflexion carries a
verbal reflection on failure into later attempts~\citep{reflexion}. In their
survey of these methods, \citet{pan2024survey} observe that when a model
produces its own feedback, the quality of that feedback is constrained by the
model's limitations, and that recent work therefore draws feedback from external
tools such as trained models, code interpreters and search engines.
Two critical studies narrow the claim. \citet{huang2024cannot} show that,
without an oracle deciding when to stop, intrinsic self-correction does not
improve reasoning and often degrades it. They also caution that feedback
messages can carry task information the original prompt omitted, which makes a
feedback effect partly a prompt-completeness effect. \citet{kamoi2024} classify
correction driven by an external, non-LLM verifier as a fair setting distinct
from intrinsic self-correction. They identify feedback generation as the
bottleneck and recommend measuring feedback quality directly. Our supervisor is
such an external verifier. Its feedback is generated deterministically and is
correct by construction about what it measures.

Whether feedback \emph{content} matters, as distinct from the signal that a
revision is needed, is where the literature divides. \citet{olausson2024selfrepair}
establish two norms for repair claims. A repair must be compared with
independent resampling at an equal budget. And feedback quality is binding:
substituting human-written for model-written explanations raised successful
repairs from 33.3\% to 52.6\%, on a deliberately easy sample of 40 failing
programs. \citet{tyen2024} separate finding an error from correcting it. Models
locate the first error in a reasoning trace poorly, but given its location, a
simple backtracking procedure corrected 18 to 44 percentage points of
previously incorrect traces, and supplying a random location roughly halved
the gain. In self-debugging, \citet{chen2024selfdebug} find that their least
informative feedback condition retains much of the benefit of the most
informative ones, because execution still decides whether to revise. Removing
that trigger can push accuracy below the no-repair baseline. Two studies
manipulate presentation rather than content. \citet{tsui2026scbench} injects the
same error either into the user's turn or into the model's own, and finds that
models correct the external error far more often. Its self-correction blind
spot, one minus the ratio of the internal to the external correction rate,
averages 64.5\% across 14 models. Appending a contentless continuation cue
(``Wait'') reduces that blind spot by 89.3\%. \citet{chen2026rolerelabel}, in
an unrefereed preprint, hold an erroneous claim byte-identical while changing
only the chat role that carries it. Explicit correction rises by 23 to 93
points, significantly in 10 of 12
model--domain combinations. Two code-repair studies shape the single-retry
design: a second editing round can degrade results~\citep{zhang2023selfedit}, and
most of the gain from iterated repair arrives in the first round or
two~\citep{arimbur2026}. \citet{arimbur2026} also documents run-to-run variation
in greedy decoding on hosted inference APIs, a threat this study measures
directly.

None of these studies holds the defect fixed and varies only what the failing
stage is told, under a deterministic external checker, with the placement of
the defective output controlled. Arms \arm{A1} and \arm{A2} of our design are
exactly that manipulation. Our placement factor reproduces the user-turn versus
own-turn contrast of \citet{tsui2026scbench} and \citet{chen2026rolerelabel},
and adds the placement that production retries actually use, in which the
defective output is absent.

\subsection{Fault injection and failure analysis in multi-agent LLM systems}
\label{sec:related-fi}

Injecting faults to observe how a system degrades has a long engineering
lineage, including chaos engineering of production services~\citep{basiri2016chaos}.
Its recent application to LLM agents has developed along two tracks.
\emph{Failure analysis} catalogues what goes wrong. \citet{cemri2025mast}
derive a taxonomy of 14 failure modes in three categories from expert annotation
of multi-agent traces, with inadequate task verification as one of the three
principal categories. They observe that verifiers often perform only
superficial checks, and that the presence of a failure mode is only loosely
coupled with whether the task succeeds. \emph{Fault injection} measures how
damage propagates. \citet{huang2025resilience} corrupt agents' profiles or
messages and find that linear structures degrade most and hierarchical ones
least. Corrupting a profile also degrades performance more than corrupting
messages, so the two injection methods are not interchangeable.
MAS-FIRE~\citep{jia2026masfire}, a preprint, injects 15 types of fault through
prompts, intercepted outputs and message routing. It separates whether a system
detects a fault, whether it resolves it locally, and whether the task still
succeeds. AgentChaos~\citep{tan2026agentchaos} injects 65 configurations of
crash, omission and value faults at the boundary between agent and LLM API.
It discards tasks in which the configured fault never fired, a trigger
verification our design adopts. It finds that injection position matters most
for pipeline architectures, and that silent failures, plausible outputs with no
warning, account for 22\% to 66\% of failures depending on the system.
A third line of work, failure attribution, asks which agent and which step
caused a failed run. Who\&When~\citep{zhang2025whowhen} is a benchmark of
hand-annotated failure logs from 127 multi-agent systems; the best of the three
attribution methods evaluated on it found the decisive step in 14.2\% of cases.
AgenTracer~\citep{zhang2026agentracer} labels runs automatically to train an
attribution model. Failed runs are labelled by counterfactual replay, which
corrects one action at a time and re-runs the rest; successful runs become
failures through an injected fault, whose step is then known by construction.
Attribution locates a fault after a run has failed. Our supervisor locates a
defect while the pipeline runs, and our injections measure whether the retry
then removes it.

This work measures tolerance, propagation and diagnosis. None of it measures
whether a supervisor's feedback gets a defect \emph{repaired}. None of it checks
whether the injected faults behave like the faults the system produces on its
own.

\subsection{Synthetic faults as proxies for real ones}
\label{sec:related-proxy}

Software engineering asked the second question of mutation analysis long
before LLM pipelines existed. Testing studies commonly used hand-seeded faults
and mutants in place of real bugs, which are hard to extract and reproduce.
Defects4J~\citep{just2014defects4j} was built to make real faults available
for controlled studies. \citet{just2014mutants} test whether artificial
mutants can stand in for real faults. Across 357 real faults in five Java
programs, 73\% were coupled to mutants from commonly used operators, and a
test suite's mutant detection correlated with its real-fault detection beyond
what coverage explains. Their coupling concerns \emph{detection}. They note that
it is unclear whether fault-localisation and program-repair techniques
evaluated on mutants perform equally well on real faults. Our design faces the
same question for repair: whether a repair rate measured on injected defects
says anything about the defects a pipeline produces. We answer it empirically.
The 13 defects the production supervisor caught are replayed through the same
retry. Before the study, blinded raters compared real and injected defects, and
both results are reported beside the injected-defect rates
(Sections~\ref{sec:design-operators} and~\ref{sec:results-rq3}).
